\begin{table}[H]
\centering
\caption{\label{tab:sample-flow}Analysis sample flow}
\centering
\begin{threeparttable}
\begin{tabular}[t]{>{\raggedright\arraybackslash}p{3.2cm}>{\raggedright\arraybackslash}p{6.6cm}>{\raggedleft\arraybackslash}p{1.2cm}}
\toprule
Sample & Stage & N\\
\midrule
FFRISP & Source export (starting N) & 902\\
FFRISP & After valid panel-ID restriction & 898\\
FFRISP & Final analytic sample (nonmissing outcome) & 887\\
Luth & Source export (starting N) & 2,330\\
Luth & After attention-check restriction & 2,017\\
Luth & After valid-instrument restriction & 1,777\\
Luth & Final analytic sample (nonmissing outcome) & 1,777\\
Luth & Sensitivity: all completed, valid outcome & 2,048\\
Stanford discovery & Source export (starting N) & 1,647\\
Stanford discovery & Final analytic sample (nonmissing first choice) & 1,645\\
Stanford replication & Source export (starting N) & 1,442\\
Stanford replication & Final analytic sample (nonmissing first choice) & 1,440\\
UCSB replication & Source export (starting N) & 1,500\\
UCSB replication & Final analytic sample (nonmissing choice) & 1,500\\
\bottomrule
\end{tabular}
\begin{tablenotes}[para]
\item Note: FFRISP's numeric-identifier step removes four test or nonpanel records. The Luth valid-pair step removes the defective natural-frequency pair. The final rows identify the primary analytic sample for each survey. The Luth sensitivity row retains completed respondents with a valid outcome without the attention screen. The 2018 rows show source-to-analysis attrition separately for each of the three study exports.
\end{tablenotes}
\end{threeparttable}
\end{table}
